% =====================================================================
%  3.6 Modellering — ekstra utfordring
%  Alle data er tenkte data.
% =====================================================================

\begin{infoboks}{Velg metode selv}
Arket bygger på 3.6 Modellering. Bruk regresjon i et digitalt verktøy der oppgaven krever det.
Skriv svar med enhet og med forbehold («ifølge modellen»). Alle data er tenkte data.
\end{infoboks}

\nivaahode{2}

\oppg{U1}{2}{\uh}
Hvilken modelltype passer? Skriv også et funksjonsuttrykk. Begrunn.
\begin{tasks}(1)
  \task Vannstanden i en tank er 30 cm og stiger 4 cm hver time.
  \task En bil har verdi 1\,000 (tusen kr) og faller $12\,\%$ i verdi hvert år.
  \task Arealet av et kvadrat når sidelengden $s$ varierer.
  \task En bakteriekultur starter med 3 enheter og dobles hver time.
\end{tasks}

\oppg{U2}{2}{\mh}
Tenkte data: vekten $W$ (gram) til en fisk med lengde $L$ (cm).

\smallskip
\begin{tabular}{@{}c|ccccc@{}}
$L$ & 10 & 15 & 20 & 30 & 40\\ \hline
$W$ & 12 & 40 & 95 & 320 & 760\\
\end{tabular}

\begin{tasks}(1)
  \task Regn ut $\dfrac{W(20)}{W(10)}$ og $\dfrac{W(40)}{W(20)}$. Hva forteller det?
  \task Bruk regresjon og skriv modellen.
  \task Hva skjer med vekten når lengden dobles?
  \task Anslå vekten til en fisk på 60 cm. Hvor sikkert er svaret?
\end{tasks}

\oppg{U3}{2}{\mh}
Tenkt avisutdrag: «Appen Handlehjelp slo an. 1. januar 2023 hadde den 50\,000 brukere. Ett år senere var tallet 90\,000, og 1. januar 2025 hadde 162\,000 lastet den ned.
– Med denne veksten vil alle i Norge bruke appen innen utgangen av 2030, sier gründeren.» Det bor omtrent 5,5 millioner mennesker i Norge.
\begin{tasks}(1)
  \task Sett tallene i en tabell. La $x$ være år etter 2023.
  \task Velg modelltype og lag en modell.
  \task Hvor mange brukere gir modellen 1. januar 2031?
  \task Er påstanden rimelig? Nevn to ting du måtte ha visst for å vurdere den.
\end{tasks}

\oppg{U4}{2}{\prog}
Programmet finner den rette linja $y = ax + b$ med minste kvadraters metode:
\[ a = \frac{\sum (x_i - \bar x)(y_i - \bar y)}{\sum (x_i - \bar x)^2}, \qquad b = \bar y - a\,\bar x \]
Tenkte data: $x$ er temperatur ($^\circ$C) og $y$ er salg av iskrem (tusen kr).
\begin{kode}
x = [14, 17, 20, 23, 26, 29]
y = [22, 31, 36, 47, 52, 63]
n = len(x)
x_snitt = sum(x) / n
y_snitt = sum(y) / n
teller = 0
nevner = 0
for i in range(n):
    teller = teller + (x[i] - x_snitt) * (y[i] - y_snitt)
    nevner = nevner + ______
a = teller / nevner
b = ______
print(a, b)
\end{kode}
\begin{tasks}(1)
  \task Fyll ut de to hullene.
  \task Kjør programmet. Sammenlign $a$ og $b$ med lineær regresjon i et digitalt verktøy.
  \task Tolk $a$ og $b$ med enhet. Hva er rart med $b$?
\end{tasks}

\nivaahode{3}

\oppg{U5}{3}{\mh}
Lag selv fem punkter med $x = 1, 2, 4, 8, 16$ der en potensmodell $y = a \cdot x^b$ passer mye bedre enn en eksponentiell modell.
Bruk regresjon for å sjekke. Forklar hvordan du valgte $y$-verdiene.

\oppg{U6}{3}{\uh}
Regresjonslinja har $b = \bar y - a\,\bar x$.
\begin{tasks}(1)
  \task Vis at linja alltid går gjennom punktet $(\bar x, \bar y)$.
  \task En regresjonslinje har stigningstall $3{,}0$. Gjennomsnittet av $x$-verdiene er 4, og gjennomsnittet av $y$-verdiene er 20. Finn $b$.
\end{tasks}

\oppg{U7}{3}{\mh}
Tenkte data: bruktpris $y$ (tusen kr) for en bilmodell etter alder $x$ (år).

\smallskip
\begin{tabular}{@{}c|cccccccc@{}}
$x$ & 1 & 2 & 3 & 4 & 5 & 6 & 7 & 8\\ \hline
$y$ & 310 & 262 & 224 & 193 & 170 & 153 & 142 & 138\\
\end{tabular}

\begin{tasks}(1)
  \task Bruk andregradsregresjon. Finn modellen og $R^2$.
  \task Hva sier modellen om prisen på en 15 år gammel bil?
  \task Prøv eksponentiell regresjon. Sammenlign $R^2$ og prognosen for 15 år. Hvilken modell bør du velge, og hvorfor?
\end{tasks}

\begin{oppsummering}
Velg modelltype ut fra tabellen og situasjonen, ikke bare ut fra $R^2$.
Bruk modellen innenfor gyldighetsområdet, og skriv «ifølge modellen» om prognoser.
Tolk alltid konstantene med enhet.
\end{oppsummering}
